% ECONOMICS_PROBLEM: {"id": "I26-607", "title": "Persistence of early-childhood benefits", "jel_code": "I26", "source_id": "OP-0450"}
% !TeX program = xelatex
\documentclass[11pt,a4paper]{article}
\usepackage[margin=23mm]{geometry}
\usepackage{fontspec}
\setmainfont{Latin Modern Roman}
\usepackage{amsmath,amssymb,mathtools}
\usepackage{xcolor,enumitem,booktabs,longtable,array,xurl,hyperref}
\definecolor{muted}{HTML}{53636C}
\hypersetup{colorlinks=true,urlcolor=blue,linkcolor=blue}
\setlength{\parindent}{0pt}
\setlength{\parskip}{5pt}
\newcommand{\R}{\mathbb R}
\newcommand{\E}{\mathbb E}
\newcommand{\N}{\mathbb N}
\newcommand{\ind}{\mathbf 1}
\newcommand{\Var}{\operatorname{Var}}
\newcommand{\Cov}{\operatorname{Cov}}
\newcommand{\supp}{\operatorname{supp}}
\DeclareMathOperator*{\argmin}{arg\,min}
\DeclareMathOperator*{\argmax}{arg\,max}
\newcommand{\entrymeta}[3]{{\sffamily\small\color{muted}\textbf{\texttt{#1}}\quad|\quad #2\par #3\par}}
\newcommand{\Problemref}[1]{\texttt{#1}}
\newcommand{\JELref}[2]{\texttt{#2} (JEL \texttt{#1})}
\begin{document}
\section*{Persistence of early-childhood benefits}
\textbf{Problem ID:} \texttt{I26-607}\quad \textbf{JEL:} \texttt{I26}\par
% BEGIN ECONOMICS STATEMENT
\label{problem:OP-0450}
{\sffamily\small\color{muted}Original subject group: Education, families and demography\par}
\label{definitions:problem:30007732}
\textbf{Audit classification:} Published research agenda. Checked source identity and appropriate agenda/background support; expanded intervention, units, population, definedness and causal restrictions. Exact targets and variants are editorial.
\subsection*{Definitions and precise research target}
\textit{Editorial formalization.} Consider children eligible for a public preschool program in one country, followed from age four to age twenty-five. Let \(p\in\{0,1\}\) denote a randomized preschool offer and \(f\in\{0,1\}\) a separately randomized improvement in primary-school instruction beginning at age six. Let \(S_{ia}(p,f)\) be standardized skills at age \(a\), and \(Y_i(p,f)\) cumulative real earnings through age twenty-five. Define
\[\tau_a(f)=E[S_{ia}(1,f)-S_{ia}(0,f)],\qquad\tau_Y(f)=E[Y_i(1,f)-Y_i(0,f)].\]
The expectation covers baseline eligible children, including those who decline preschool, with outcome scales fixed independently of treatment. Identify the skill-response trajectory and earnings effects under factorial assignment, consistency, and a specified school-level spillover structure. Determine whether later instructional quality changes persistence of early gains, and whether convergence of measured test scores predicts convergence of later economic outcomes. Distinguish measurement artifacts, replacement of skills, and compensating parental investment where additional data permit. Report attrition and administrative-linkage uncertainty.

Fix outcome ages \(a\in\{5,6,10,15,20,25\}\); use a common independently calibrated scale at each age, declaring which skills are comparable over time. Earnings \(Y_i\) are \(\sum_{a=16}^{25}\beta^{a-4}y_{ia}\) at common real prices, including zeros for nonemployment. A factorial design randomizes both preschool offer \(p\) and follow-on school assignment \(f\), with positive probabilities in all four arms; separate school clusters may require two-stage assignment and explicit interference rules. \(\tau_a(1)-\tau_a(0)\) is an offer interaction on the specified skill scale. Test-score fadeout means this measured contrast shrinks, not disappearance of every latent skill or earnings gain. Follow-on schooling choice and parental response are outcomes unless separately randomized. Specify mortality/attrition and administrative linkage rather than conditioning on eventual observation. All expectations use a stated baseline population and probability law. Identification means every admissible data-generating model with the same observed-data law yields the same displayed target; otherwise report the identified set. Exact equations, horizons and assumptions are editorial, rather than source-given canonical conjectures.
\subsection*{Variants and assumption changes}
\begin{enumerate}
\item \textbf{Editorial latent-skill variant.} Use a longitudinal multiple-skill measurement model with invariant item parameters. Test invariance before interpreting score convergence as skill convergence.
\item \textbf{Editorial sequential-investment variant.} Randomize later schooling after observing preschool outcomes with known probabilities. Compare adaptive policies with fixed factorial policies under sequential randomization.
\end{enumerate}
\subsection*{References and source audit}
\textbf{1.} Explicit agenda on long-run curricula, follow-on investments, and combined preschool/parenting programs. Locator: BFI Working Paper 2022-58; Section 7, printed pp.91–93, particularly p.93. Full text or primary publication page inspected. URL: \url{https://bfi.uchicago.edu/wp-content/uploads/2022/05/BFI_WP_2022-58.pdf}.
\begin{enumerate}\raggedright
\item\hypertarget{definitions:source:30007732:1}{} Greg J. Duncan; Ariel Kalil; Magne Mogstad; Mari Rege. 2022. \emph{Investing in Early Childhood Development in Preschool and at Home}. BFI Working Paper 2022-58; Section 7, printed pp.91–93, particularly p.93.
\url{https://bfi.uchicago.edu/wp-content/uploads/2022/05/BFI_WP_2022-58.pdf}.
DOI: \url{https://doi.org/10.3386/w29985}.
\end{enumerate}

\subsection*{Common conventions: Source mathematical conventions}

These conventions apply to each entry unless explicitly replaced.

\textbf{Models and identification.} A primitive model \(m\) specifies all probability laws, feasible choices, information, equilibrium and measurement rules used by its target. The admissible class \(\mathcal M\) is fixed independently of outcomes. If \(P_O(m)\) is its observable probability law and \(\tau(m)\) the target, its identified set at observable law \(P\) is
\[
 \mathcal I_\tau(P)=\{\tau(m):m\in\mathcal M,\ P_O(m)=P\}.
\]
Point identification means that this set is a singleton. An empty set signals incompatibility of data and assumptions. Coordinate bounds need not describe the joint set. Bounds are sharp only if every retained target is attainable or arbitrarily approachable by compatible models. Finite bounds require explicit restrictions on unobserved quantities; unrestricted confounding, hidden wealth, extrapolation or arbitrary equilibrium selection can yield uninformative sets.

\textbf{Causal interventions.} A potential outcome \(Y_i(p)\) is the outcome under a complete policy regime \(p\), including timing, financing, information and market spillovers. The target population and units are fixed before treatment. With interference, use the complete assignment vector or a justified exposure mapping; individual no-interference notation alone cannot identify an economy-wide intervention. A difference of mean potential outcomes does not require the cross-world joint distribution. Conditioning on post-treatment migration, survival, enrollment or employment changes the estimand and needs separate justification.

\textbf{Statistical statements.} An estimator, confidence set, forecast or test specifies a sampling law, model class and loss/coverage criterion. Uniform \((1-\alpha)\)-coverage means \(\inf_{m\in\mathcal M}\Pr_m\{\tau(m)\in C_n\}\geq1-\alpha\), or an explicitly asymptotic version. The whole parameter space trivially has coverage, so informativeness requires a stated width, risk or power objective. A forecasting improvement is relative to a named benchmark, information set and evaluation population.

\textbf{Equilibrium and optimization.} An equilibrium comprises feasible plans, optimality, beliefs where needed, budget constraints and all market/resource clearing. The primitives or a stated benchmark define each correspondence \(\mathcal E(p;m)\). Nonemptiness is assumed or proved, never inferred just from compact policy sets. A selected-equilibrium target uses a declared selection rule; a robust target quantifies over all permitted equilibria. Use suprema or infima unless attainment is established. An \(\varepsilon\)-optimal policy is feasible and within \(\varepsilon\) of the finite optimum value. Report an empty feasible set separately.

\textbf{Welfare and accounting.} Normative utility, interpersonal normalization, social weights, numeraire, discounting and the use of public receipts are primitives. Behavioral utility need not coincide with the welfare criterion. Household welfare includes ownership income and rents once; adding profits again to compensating variation generally double-counts them. A monetary equivalent is defined by an explicit transfer and price convention, with monotonicity and range conditions. Average effects, mechanism contributions, transfers and welfare losses are different targets. Infinite sums require convergence and dynamic feasible plans require terminal or no-Ponzi conditions.

\textbf{Source versions.} A working-paper date, revision date and journal publication year are distinguished. A DOI for a journal version is not automatically the DOI of an earlier manuscript. Locators refer to the stated version and printed pages where specified. A metadata-only check does not verify a claim about a full-text conclusion. Every entry's source subsection narrows the provenance claim accordingly.
% END ECONOMICS STATEMENT
\end{document}
